\section{OUTLINE}
\label{sec:outline}

To present these contributions and delineate the current progress of this research, the remainder of this qualification document is structured as follows: Chapter \ref{chap:observability} reviews the historical evolution of observability from traditional control theory to modern distributed systems, defining the core telemetry pillars and detailing the analytical shift required to transition from conventional monitoring to true system observability. Next, Chapter \ref{chap:related-work} presents a Rapid Review of the state-of-the-art literature, systematically mapping existing solutions to highlight current architectural gaps and introducing the Observability Maturity Model (OMM) to categorize the evolutionary stages of telemetry management. Chapter \ref{chap:cloud-computing} then provides the theoretical background on Cloud-Native computing and the Kubernetes orchestration model, also covering the varying levels of instrumentation, Infrastructure as Code (IaC) principles, and Chaos Engineering as a resilience validation mechanism. Following this, Chapter \ref{chap:methodology} details the methodology of the study, outlining the Design Science Research (DSR) framework, formulating the research hypotheses, and defining the experimental design alongside the specific metrics selected to evaluate the proposed solution. Chapter \ref{chap:experiment-implementation} describes the technical instantiation of the experimental testbed and the PANOPTES reference architecture, detailing the virtualization layer, the IaC deployment strategy, the configuration of the observability components, and the fault injection setup. Chapter \ref{chap:results-and-discussion} unveils the discoveries using RCA algorithmic strategy to validade research hypothesis through each scenarios presented in this work.  Finally, Chapter \ref{chap:conclusion} summarizes the architectural decisions consolidated thus far, assesses the current state of the research, and outlines the scheduled milestones and future work required leading up to the final thesis defense.